% from ega2-1-fr.tex
\oldpage[II]{120}

$L'$ (resp. $K'$) de $L$ (resp. $K$) dans $L_1$ (resp. $K_1$), et il existe
un morphisme canonique (nécessairement dominant) $f':X'\to Y'$ rendant
commutatif le diagramme (\hyperref[II.6.3.5.1-fr]{6.3.5.1}).

Le cas le plus important est celui où on prend $L_1=L$, $K_1$ étant donc une
extension de $K$ contenue dans $L$, et où on suppose $X$ intègre et
\emph{normal}, donc $X'=X$. Ce qui précède montre alors que lorsque $X$ est
normal, et que $Y'$ est la fermeture intégrale de $Y$ relative à un corps
$K_1\subset L=R(X)$, tout morphisme dominant $f:X\to Y$ se factorise en
\[
f:X\xrightarrow{f'}Y'\to Y
\]
où $f'$ est dominant~; d'ailleurs, lorsque le monomorphisme $K_1\to L$ est
donné, $f'$ est nécessairement unique, comme on le voit en se ramenant au cas
où $X$ et $Y$ sont affines. On voit donc que pour la donnée de $Y$, $L$ et
d'un $K$-monomorphisme $K_1\to L$, la fermeture intégrale $Y'$ de $Y$
relativement à $K_1$ est solution d'un \emph{problème universel}.
\end{env}

\begin{remark}[6.3.10]
\label{II.6.3.10-fr}
Reprenons les hypothèses de (\hyperref[II.6.3.6-fr]{6.3.6}), en supposant de
plus que chacune des algèbres $A_i$ soit de \emph{rang fini} sur $k(\xi_i)$
(ce qui entraîne que $A_i$ est composée directe d'un nombre fini de corps)~;
on peut affirmer dans certains cas que le morphisme structural $X'\to X$ est
non seulement entier, mais même \emph{fini}. Bornons-nous au cas où $X$ est
\emph{réduit}~; la question étant locale sur $X$, on peut en outre supposer
que $X$ est affine d'anneau $C$, et que $C$ n'a qu'un nombre fini d'idéaux
minimaux $\mathfrak{p}_i$ ($1\leq i\leq r$) tels que les
$C_i=C/\mathfrak{p}_i$ soient intègres~; $X'$ sera alors fini sur $X$ si la
fermeture intégrale de chaque $C_i$ dans toute extension de degré fini de son
corps des fractions est un $C$-module de type fini
(\hyperref[II.6.3.4-fr]{6.3.4}). On sait que cette condition est toujours
vérifiée lorsque $C$ est une \emph{algèbre de type fini sur un corps}
(\cite[t.~I, p.~267, th.~9]{I-13}), ou \emph{sur $\mathbf{Z}$}
(\cite[t.~I, p.~93, th.~3]{I-9}), ou \emph{sur un anneau local noethérien
complet} ([25], p.~298). On en conclut que $X'\to X$ sera un morphisme
fini lorsque $X$ est un schéma de \emph{type fini sur un corps}, ou sur
$\mathbf{Z}$, ou sur un anneau local noethérien complet.
\end{remark}

\subsection{Déterminant d'un endomorphisme de $\mathcal{O}_X$-Module.}
\label{subsection:II.6.4-fr}

\begin{env}[6.4.1]
\label{II.6.4.1-fr}
Soient $A$ un anneau (commutatif), $E$ un $A$-module libre de rang $n$, $u$
un endomorphisme de $E$~; rappelons que pour définir le \emph{polynôme
caractéristique} de $u$, on considère l'endomorphisme $u\otimes1$ du
$A[T]$-module libre de rang $n$, $E\otimes_AA[T]$ ($T$ indéterminée), et on
pose
\[
\label{II.6.4.1.1-fr}
P(u,T)=\det(T\cdot I-(u\otimes1))
\tag{6.4.1.1}
\]
($I$ automorphisme identique de $E\otimes_AA[T]$). On a
\[
\label{II.6.4.1.2-fr}
P(u,T)=T^n-\sigma_1(u)T^{n-1}+\ldots+(-1)^n\sigma_n(u)
\tag{6.4.1.2}
\]
où $\sigma_i(u)$ est un élément de $A$, égal à un polynôme homogène de degré
$i$ (à coefficients entiers) par rapport aux éléments de la matrice de $u$
relative à une base quelconque de $E$~; on dit que les $\sigma_i(u)$ sont les
\emph{fonctions symétriques élémentaires de $u$}, et on a en particulier
$\sigma_1(u)=\operatorname{Tr}u$ et $\sigma_n(u)=\det u$. Rappelons qu'en
vertu du th. de Hamilton-Cayley, on a
\[
\label{II.6.4.1.3-fr}
P(u,u)=u^n-\sigma_1(u)u^{n-1}+\ldots+(-1)^n\sigma_n(u)=0
\tag{6.4.1.3}
\]

